% ============================================================
%  Modul 10 — Implementasi Mutex dan Locks
% ============================================================

\chapter{Implementasi Mutex dan Locks}
\label{chap:modul-10}

% --- 1. Tujuan dan Output Praktikum ---
\subsection{Tujuan Praktikum}
Setelah menyelesaikan praktikum ini, mahasiswa diharapkan mampu:
\begin{enumerate}
    \item Memahami cara kerja \textit{semaphore} sebagai mekanisme dasar sinkronisasi antar \textit{thread}.
    \item Memahami konsep \textit{mutex} (\textit{Mutual Exclusion Lock}) dan perbedaannya dengan \textit{semaphore} secara semantik maupun struktural.
    \item Mengimplementasikan mekanisme penahanan (\textit{wait}) dan pelepasan (\textit{signal}) sumber daya, serta sistem penguncian eksklusif pada tingkat kernel.
    \item Memverifikasi efektivitas mekanisme sinkronisasi yang dibangun dalam mencegah dan menghilangkan \textit{race condition} yang diamati pada modul sebelumnya.
\end{enumerate}

\subsection{Output Praktikum}
Pada akhir praktikum ini, mahasiswa diharapkan menghasilkan:
\begin{itemize}
    \item Implementasi algoritma sinkronisasi kernel yang valid dan mampu meregulasi antrian \textit{thread} dengan tepat.
    \item Simulasi sistem operasi yang mendemonstrasikan penyelesaian skenario \textit{Producer-Consumer} dengan eksekusi yang tersinkronisasi, rapi, dan konsisten.
    \item Laporan analisis yang menjabarkan bagaimana struktur \textit{semaphore} dan \textit{lock} bekerja sama secara harmonis untuk melindungi area \textit{critical section} dari interupsi paksa.
\end{itemize}

% --- 2. Dasar Teori ---
\newpage
\section{Dasar Teori}

\subsection{Konsep Semaphore}

\textit{Semaphore} adalah sebuah struktur data abstrak yang digunakan untuk mengkoordinasikan akses beberapa \textit{thread} atau proses terhadap sumber daya (\textit{shared resource}) yang jumlahnya terbatas.
Konsep ini pertama kali diperkenalkan oleh ilmuwan komputer Edsger W. Dijkstra pada tahun 1965.

Secara konseptual, \textit{semaphore} bekerja seperti "keranjang berisi token". Jika sebuah \textit{thread} ingin melakukan pekerjaan, ia harus mengambil satu token dari keranjang. Jika keranjang kosong,
\textit{thread} tersebut harus menunggu (mengantri) sampai ada token yang dikembalikan oleh \textit{thread} lain.

Di dalam sistem operasi, sebuah \textit{semaphore} terdiri dari dua komponen utama:
\begin{enumerate}
    \item Sebuah variabel bilangan bulat (integer) non-negatif yang mewakili "jumlah token" atau kapasitas sumber daya.
    \item Sebuah antrian (\textit{waiting queue}) untuk menampung \textit{thread} yang sedang menunggu giliran.
\end{enumerate}

Semaphore dioperasikan secara mutlak melalui dua operasi atomik (tidak dapat diinterupsi):
\begin{itemize}
    \item \textbf{Operasi \texttt{Wait(S)} (P atau \textit{Proberen})}: Mengurangi nilai variabel semaphore. Jika nilainya menjadi $0$ (atau kurang), maka \textit{thread} yang mengeksekusi operasi ini akan diubah
        statusnya menjadi \textit{Blocked} / \textit{Asleep} dan dimasukkan ke dalam antrian. Pada sistem OS, operasi ini dinamakan \texttt{sema\_down()}.
    \item \textbf{Operasi \texttt{Signal(S)} (V atau \textit{Verhogen})}: Menambah nilai variabel semaphore. Jika ada \textit{thread} yang sedang mengantri (tertidur), operasi ini akan membangunkan (\textit{Wake up})
        salah satu \textit{thread} tersebut agar siap dieksekusi kembali oleh CPU. Pada sistem OS, operasi ini dinamakan \texttt{sema\_up()}.
\end{itemize}

Berdasarkan kapasitasnya, semaphore dibagi menjadi dua:
\begin{enumerate}
    \item \textbf{Counting Semaphore}: Nilainya bisa berapa saja ($\geq 0$). Sangat berguna untuk mengontrol akses ke *resource* yang memiliki batasan jumlah, seperti jumlah tempat parkir kosong, atau sisa kapasitas \textit{buffer}.
    \item \textbf{Binary Semaphore}: Nilainya hanya dibatasi pada angka $0$ dan $1$. Sering digunakan untuk menjamin \textit{Mutual Exclusion}.
\end{enumerate}

\subsection{Mutex (Mutual Exclusion Lock)}

\textit{Mutex} atau \textit{Lock} adalah mekanisme sinkronisasi yang secara khusus dirancang untuk menjamin \textit{Mutual Exclusion} pada area \textit{Critical Section}. 

Analogi sederhana dari \textit{Mutex} adalah kunci kamar mandi umum. 
\begin{itemize}
    \item Jika ada orang (Thread A) masuk ke kamar mandi (\textit{Critical Section}), ia akan mengunci pintunya (\textit{Lock}).
    \item Orang lain (Thread B) yang ingin masuk harus mengantri di luar sampai pintu dibuka.
    \item Pintu hanya bisa dibuka dari dalam oleh orang yang sama yang menguncinya tadi (Thread A).
\end{itemize}

\subsubsection{Konsep Ownership (Kepemilikan)}
Ini adalah perbedaan mendasar antara Mutex dan Semaphore. Mutex memiliki konsep \textbf{Ownership} (Kepemilikan). Sebuah Mutex akan mengingat identitas \textit{thread} mana yang berhasil menguncinya. Sistem operasi yang baik
akan menolak operasi \textit{Unlock} jika fungsi tersebut dipanggil oleh \textit{thread} yang bukan pemiliknya. 

Di dalam level implementasi (seperti pada OS), Mutex sering kali dibangun di atas \textit{Binary Semaphore} dengan nilai awal $1$, namun ditambah dengan variabel \texttt{holder} untuk mencatat siapa pemegang kunci tersebut.

\begin{table}[h]
\centering
\caption{Perbedaan Konseptual Semaphore dan Mutex}
\begin{tabulary}{\textwidth}{|L|L|L|}
\hline
\textbf{Aspek} & \textbf{Semaphore} & \textbf{Mutex (Lock)} \\ \hline
Analogi & Keranjang berisi token bebas. & Kunci gembok eksklusif. \\ \hline
Konsep Nilai & Bilangan bulat $\geq 0$ (tergantung kapasitas). & Hanya $0$ (Terkunci) atau $1$ (Bebas). \\ \hline
Kepemilikan (\textit{Ownership}) & Tidak ada. Thread A bisa melakukan \texttt{Wait}, dan Thread B bisa melakukan \texttt{Signal}. & Ada. Hanya thread pemegang kunci yang berhak melakukan \textit{Release/Unlock}. \\ \hline
Kegunaan Utama & Sinkronisasi urutan kejadian dan manajemen jumlah sumber daya. & Melindungi \textit{Critical Section} dari modifikasi serentak. \\ \hline
\end{tabulary}
\end{table}

\subsection{Pendekatan Logika Sinkronisasi Producer-Consumer}

Pada Modul 9, kita melihat bahwa masalah \textit{Producer-Consumer} (\textit{Bounded-Buffer}) menyebabkan anomali ketika \textit{buffer} diakses tanpa aturan. Untuk menyelesaikannya kita membutuhkan gabungan dari Semaphore dan Mutex.

Logika penyelesaiannya membutuhkan 3 instrumen:
\begin{enumerate}
    \item \textbf{\texttt{mutex\_lock}}: Untuk melindungi variabel \textit{buffer} agar tidak ditulis/dibaca bersamaan.
    \item \textbf{\texttt{sema\_empty}}: \textit{Counting Semaphore} untuk menghitung jumlah slot kosong di dalam \textit{buffer}.
    \item \textbf{\texttt{sema\_full}}: \textit{Counting Semaphore} untuk menghitung jumlah slot terisi di dalam \textit{buffer}.
\end{enumerate}

\noindent \textbf{Alur Logika Producer:}
\begin{itemize}
    \item Sebelum memproduksi barang, Producer harus \texttt{Wait(sema\_empty)}. Jika kapasitas penuh (kosong = 0), ia harus tidur.
    \item Jika ada slot kosong, ia mengunci area dengan \texttt{Lock(mutex)}.
    \item Setelah barang ditaruh, ia membuka \texttt{Unlock(mutex)}.
    \item Terakhir, ia memberi sinyal \texttt{Signal(sema\_full)} untuk memberitahu Consumer bahwa ada barang baru.
\end{itemize}

\noindent \textbf{Alur Logika Consumer:}
\begin{itemize}
    \item Sebelum membaca barang, Consumer harus \texttt{Wait(sema\_full)}. Jika tidak ada barang (terisi = 0), ia harus tidur.
    \item Jika ada barang, ia mengunci area dengan \texttt{Lock(mutex)}.
    \item Setelah barang diambil, ia membuka \texttt{Unlock(mutex)}.
    \item Terakhir, ia memberi sinyal \texttt{Signal(sema\_empty)} untuk memberitahu Producer bahwa ada slot yang baru saja kosong.
\end{itemize}

\begin{table}[h]
\caption{Tabel Ringkasan Dasar Teori Modul 10}
\label{tab:m10-dasar-teori}
\centering
\begin{tabulary}{\textwidth}{|L|L|}
\hline
\textbf{Konsep} & \textbf{Penjelasan Singkat} \\ \hline
Wait / P() / sema\_down
    & Operasi menahan/menurunkan resource. Memblokir thread jika resource habis. \\ \hline
Signal / V() / sema\_up
    & Operasi melepaskan/menambah resource. Membangunkan thread dari antrian. \\ \hline
Ownership
    & Konsep kepemilikan pada Mutex yang memastikan hanya sang pengunci yang boleh membuka kunci. \\ \hline
Thread Blocking
    & Keadaan di mana thread diberhentikan paksa oleh OS dan dikeluarkan dari penjadwalan CPU hingga event sinkronisasi terjadi. \\ \hline
\end{tabulary}
\end{table}


% --- 3. Alat dan Bahan ---
\newpage
\section{Alat dan Bahan}
\label{sec:m09-alat-bahan}

\subsection{Perangkat Keras}
\begin{itemize}
    \item PC atau Laptop dengan arsitektur prosesor basis \textit{host} x86/x64 yang mendukung kelanjutan kompilasi emulator.
    \item Memori minimal 4GB RAM ruang eksekusi.
\end{itemize}

\subsection{Perangkat Lunak (Bahan Praktikum)}
\begin{itemize}
    \item OS \textit{Host} utama (Windows dengan WSL/GitBash, Linux, atau macOS).
    \item Text Editor (misalnya Visual Studio Code) untuk menyunting berkas kode bahasa C dan \textit{Assembly}.
    \item Instalasi Docker Engine (\texttt{docker.io}) untuk menyalakan lingkungan isolasi kompilasi.
    \item Emulator Hardware x86: \texttt{qemu-system-x86} (atau instalasi QEMU setara) yang akan mengeksekusi binari hasil OS.
    \item Direktori \textit{Starter Code} OS 32-bit yang sudah disiapkan sebelumnya (berisi direktori \texttt{src} dengan kerangka lengkap \texttt{kernel.c}, \texttt{thread.h}, \texttt{thread.c}, \texttt{switch.asm}, \texttt{scheduler.h}, dsb).
\end{itemize}

% --- 4. Langkah-Langkah Praktikum ---
\newpage
\section{Langkah-Langkah Praktikum}

\subsection{Persiapan Lingkungan}
\label{sec:m10-persiapan}

Pada praktikum ini, Anda akan menuliskan implementasi fungsi sinkronisasi di dalam \texttt{src/synch.c} dan menghubungkannya dengan simulasi di \texttt{src/kernel.c}.

\begin{enumerate}
    \item Buka terminal, pastikan Anda berada di dalam direktori \textit{source code} OS.
    \item Bersihkan hasil \textit{build} sebelumnya agar kompilasi berikutnya menggunakan kode terbaru yang akan Anda tulis.
    \begin{lstlisting}[language=bash, caption={Bersihkan build sebelumnya}]
make clean
    \end{lstlisting}
\end{enumerate}

% -----------------------------------------------------------------
\subsection{Sesi 1: Mengimplementasikan Logika Semaphore}
\label{sec:m10-sesi-1}

Semaphore adalah fondasi utama dari mekanisme sinkronisasi di OS. Buka file \texttt{src/synch.c} menggunakan teks editor Anda dan cari bagian blok fungsi \texttt{sema\_down} dan \texttt{sema\_up}.

\begin{enumerate}
    \item \textbf{Implementasi \texttt{sema\_down()} (Operasi Wait / P)} \\
    Fungsi ini bertugas menahan \textit{thread} jika \textit{resource} sedang kosong (nilai 0). Ubah fungsi \texttt{sema\_down()} menjadi seperti berikut:

    \begin{lstlisting}[language=C, caption={Kode implementasi sema\_down()}]
void sema_down(struct semaphore *sema) {
    /* 1. Selama nilai 0, masukkan thread ke antrian dan tidurkan */
    while (sema->value == 0) {
        sema->waiters[sema->waiter_count++] = thread_current();
        thread_block(); // Serahkan CPU ke thread lain
    }
    /* 2. Jika sudah lolos (nilai > 0), kurangi nilai resource */
    sema->value--;
}
    \end{lstlisting}

    \item \textbf{Implementasi \texttt{sema\_up()} (Operasi Signal / V)} \\
    Fungsi ini bertugas menambah nilai \textit{resource} dan membangunkan \textit{thread} yang sedang antri (jika ada). Ubah fungsi \texttt{sema\_up()} menjadi seperti berikut:

    \begin{lstlisting}[language=C, caption={Kode implementasi sema\_up()}]
void sema_up(struct semaphore *sema) {
    /* 1. Kembalikan/tambah nilai resource */
    sema->value++;
    
    /* 2. Jika ada thread yang menunggu, bangunkan yang pertama (FIFO) */
    if (sema->waiter_count > 0) {
        struct thread *t = sema->waiters[0];
        
        // Geser sisa antrian ke depan (menghapus index 0)
        for (uint32_t i = 0; i < sema->waiter_count - 1; i++) {
            sema->waiters[i] = sema->waiters[i + 1];
        }
        sema->waiter_count--;
        
        // Bangunkan thread agar masuk ke ready queue lagi
        thread_unblock(t);
    }
}
    \end{lstlisting}
\end{enumerate}

% -----------------------------------------------------------------
\subsection{Sesi 2: Mengimplementasikan Mutex Lock}
\label{sec:m10-sesi-2}

Setelah Semaphore berfungsi, kita akan membungkusnya menjadi sebuah Mutex Lock. Masih di dalam file \texttt{src/synch.c}, cari bagian \texttt{lock\_acquire} dan \texttt{lock\_release}.

\begin{enumerate}
    \item \textbf{Implementasi \texttt{lock\_acquire()}} \\
    Gunakan \texttt{sema\_down} untuk mengunci \textit{binary semaphore} milik lock, lalu catat siapa pemiliknya.

    \begin{lstlisting}[language=C, caption={Kode implementasi lock\_acquire()}]
void lock_acquire(struct lock *lock) {
    sema_down(&lock->semaphore);     // Tunggu & ambil kunci
    lock->holder = thread_current(); // Catat pemegang kunci saat ini
}
    \end{lstlisting}

    \item \textbf{Implementasi \texttt{lock\_release()}} \\
    Hapus kepemilikan lock, lalu gunakan \texttt{sema\_up} untuk membebaskan kunci agar bisa dipakai \textit{thread} lain.

    \begin{lstlisting}[language=C, caption={Kode implementasi lock\_release()}]
void lock_release(struct lock *lock) {
    lock->holder = NULL;             // Hapus identitas pemegang kunci
    sema_up(&lock->semaphore);       // Lepas kunci
}
    \end{lstlisting}
    \item \textbf{Simpan} file \texttt{synch.c}. Saat ini \textit{backend} sinkronisasi Anda sudah selesai.
\end{enumerate}

% -----------------------------------------------------------------
\subsection{Sesi 3: Integrasi Sinkronisasi pada Simulasi Kernel}
\label{sec:m10-sesi-3}

Mekanisme yang sudah Anda buat tidak akan berguna jika tidak dipanggil oleh sistem. Sekarang, buka file \texttt{src/kernel.c}. Kita akan memperbaiki \textit{critical section} dari skenario Producer-Consumer.

\begin{enumerate}
    \item Cari fungsi \texttt{producer(void *aux)}. Hapus tanda komentar (uncomment) pada blok TODO yang memanggil \texttt{sema} dan \texttt{lock} sehingga kodenya menjadi seperti ini:

    \begin{lstlisting}[language=C, caption={Melindungi Critical Section pada Producer}]
static void producer(void *aux) {
    (void)aux;
    int item = 0;
    while (1) {
        /* 1. TUNGGU SLOT KOSONG & KUNCI GEMBOK */
        sema_down(&sema_empty);
        lock_acquire(&buf_lock);

        /* --- CRITICAL SECTION --- */
        shared_buffer = item;
        thread_yield(); // Jebakan ini sekarang AMAN karena dilindungi lock
        vga_set_color(VGA_COLOR_BRIGHT_GREEN);
        vga_puts("[Producer] menulis item ");
        vga_putint(item++);
        vga_putchar('\n');
        /* --- END CRITICAL SECTION --- */

        /* 2. BUKA GEMBOK & BERI SINYAL ADA BARANG */
        lock_release(&buf_lock);
        sema_up(&sema_full);

        thread_yield();
    }
}
    \end{lstlisting}

    \item Lakukan hal yang sama pada fungsi \texttt{consumer(void *aux)}.

    \begin{lstlisting}[language=C, caption={Melindungi Critical Section pada Consumer}]
static void consumer(void *aux) {
    (void)aux;
    while (1) {
        /* 1. TUNGGU BARANG ADA & KUNCI GEMBOK */
        sema_down(&sema_full);
        lock_acquire(&buf_lock);

        /* --- CRITICAL SECTION --- */
        vga_set_color(VGA_COLOR_YELLOW);
        vga_puts("[Consumer] membaca item ");
        thread_yield(); // Jebakan ini sekarang AMAN karena dilindungi lock
        vga_putint(shared_buffer);
        vga_putchar('\n');
        /* --- END CRITICAL SECTION --- */

        /* 2. BUKA GEMBOK & BERI SINYAL SLOT KOSONG */
        lock_release(&buf_lock);
        sema_up(&sema_empty);

        thread_yield();
    }
}
    \end{lstlisting}
    \item \textbf{Simpan} file \texttt{kernel.c}.
\end{enumerate}

% -----------------------------------------------------------------
\subsection{Sesi 4: Kompilasi dan Verifikasi Sistem OS}
\label{sec:m10-sesi-4}

\begin{enumerate}
    \item Lakukan kompilasi sistem operasi Anda dengan \textit{flag} sinkronisasi.
    \begin{lstlisting}[language=bash, caption={Compile dan jalankan dengan sinkronisasi penuh}]
# pastikan sebelumnya sudah melakukan make clean
make DEMO=-DSYNCHRONIZATION
make run
    \end{lstlisting}

    \item \textbf{Verifikasi Hasil:} Amati output di layar QEMU. 
    \begin{itemize}
        \item \textbf{Sebelumnya (Modul 9):} Teks bertumpuk dan \textit{error} karena interupsi \texttt{thread\_yield()}.
        \item \textbf{Sekarang (Modul 10):} Setiap \texttt{"[Producer] menulis item N"} selalu diikuti secara rapi oleh \texttt{"[Consumer] membaca item N"}.
    \end{itemize}
    Meskipun CPU dipaksa berpindah \textit{thread} di tengah jalan oleh jebakan \texttt{thread\_yield()}, sistem operasi Anda kini sudah memiliki pengamanan memori.
    
    \item \textit{Screenshot} bukti hasil jalan nya untuk dilihat dengan lebih baik apakah implementasi sudah sesuai.
\end{enumerate}

\begin{figure}[h]
    \centering
    \includegraphics[width=0.85\textwidth]{Figure/10/1.jpeg}
    \caption{Output QEMU setelah implementasi sinkronisasi penuh --- Producer dan Consumer bergantian secara rapi, Mutex berhasil mengatasi jebakan Preemption.}
    \label{fig:m10-sesi4}
\end{figure}
